\honest{The Argument from Common Usage}{Eliezer Yudkowsky}{2008}

In the falling-tree dispute Albert says that people would call his recording a ``sound'': ``That's the common usage!'' Barry asks who gave Albert the right to decide what common usage is.

Most disputes, I say, reach for a dictionary because people believe words have meanings and the dictionary records them. Some even believe the dictionary sets them, maybe because a teacher once said so. But dictionary editors ``are historians.'' ``The Oxford English Dictionary may be comprehensive, but never authoritative.'' \nb{The OED calls itself ``the accepted authority on the English language'' and also ``a historical dictionary.'' It means authority as a reliable record. The post uses ``authoritative'' only for a lawgiver.} The phrase ``authoritative dictionary'' is used correctly, I say, for the IEEE's, whose members negotiate their terms. \nb{Some dictionaries are made to legislate. The French Academy's statutes of 1635 charge it with giving the language fixed rules.}

I grant the objection: shared words are a public good. It does not matter whether we say ``oto'' or ``sound,'' only that we say the same thing. Still, I can translate your usage if I know it, and people stuck on a deserted island will draw pictures in the sand.

Albert appeals to common usage only to accuse Barry of breaking it. But each knows what the other means, and both have words, ``acoustic vibrations'' and ``auditory experience,'' that describe the forest without ambiguity. So their fight cannot be about communication, the one benefit common usage gives. If only a name is at stake, coin two new words.

Often, though, a category carries hidden inferences, as when someone wants atheism counted as a religion, or moral stakes, as in the fight to class blacks and whites together as ``people.'' ``But once any empirical proposition is at stake, or any moral proposition, you can no longer appeal to common usage.'' \nb{That holds for questions about the things named. A question about what a speaker or a law meant by a word is also empirical, and common usage answers it. In \textsc{Nix v. Hedden} (1893) the Supreme Court held tomatoes to be vegetables under a tariff act because ``in the common language of the people'' they are, though botanically they are fruit.} If everyone believed that fire is the release of phlogiston, the dictionary would say so.

Definitions, I conclude, are wiser or more foolish, and treating dictionaries as the arbiters binds you to the wisdom of the past. But if you depart from it, take care that people can figure out what you are trying to swim.
